\chapter{Product Measures}

\section{Product $\sigma$-Algebras}

\begin{intuition}
    Given measurable spaces $(X, \mathcal{M})$ and $(Y, \mathcal{N})$, we want to define a natural 
    $\sigma$-algebra on $X \times Y$ that:
    \begin{enumerate}
        \item Contains ``rectangles'' $A \times B$ for measurable $A \subseteq X$ and $B \subseteq Y$
        \item Makes the canonical projections $\pi_X: X \times Y \to X$ and $\pi_Y: X \times Y \to Y$ measurable
        \item Is the \textbf{smallest} $\sigma$-algebra satisfying these properties
    \end{enumerate}
    
    It turns out these conditions are equivalent: the $\sigma$-algebra generated by rectangles 
    is precisely the smallest $\sigma$-algebra making projections measurable.
\end{intuition}

\bigskip
\begin{definition}[Canonical Projections]
    For measurable spaces $(X, \mathcal{M})$ and $(Y, \mathcal{N})$, the \textbf{canonical projections} are:
    \begin{align*}
        \pi_X: X \times Y &\to X, \quad \pi_X(x, y) = x \\
        \pi_Y: X \times Y &\to Y, \quad \pi_Y(x, y) = y
    \end{align*}
\end{definition}

\bigskip
\begin{definition}[Product $\sigma$-Algebra]
    Let $(X, \mathcal{M})$ and $(Y, \mathcal{N})$ be measurable spaces. 
    The \textbf{product $\sigma$-algebra} is:
    \[
        \mathcal{M} \otimes \mathcal{N} = \GenSigmaAlg{\{A \times B \mid A \in \mathcal{M}, B \in \mathcal{N}\}}
    \]
\end{definition}

\bigskip
\begin{theorem}[Characterization via Projections]
    The product $\sigma$-algebra is the smallest $\sigma$-algebra on $X \times Y$ making both 
    canonical projections measurable:
    \[
        \mathcal{M} \otimes \mathcal{N} = \GenSigmaAlg{\pi_X^{-1}(\mathcal{M}) \cup \pi_Y^{-1}(\mathcal{N})}
    \]
    where $\pi_X^{-1}(\mathcal{M}) = \{\pi_X^{-1}(A) \mid A \in \mathcal{M}\}$ and similarly for $\pi_Y$.
\end{theorem}

\begin{proof}
    Let $\mathcal{S} = \GenSigmaAlg{\pi_X^{-1}(\mathcal{M}) \cup \pi_Y^{-1}(\mathcal{N})}$ 
    be the smallest $\sigma$-algebra making projections measurable.
    
    \bigskip
    \textbf{Step 1: $\mathcal{M} \otimes \mathcal{N} \subseteq \mathcal{S}$.}
    
    Every rectangle $A \times B$ with $A \in \mathcal{M}$, $B \in \mathcal{N}$ satisfies:
    \[
        A \times B = \pi_X^{-1}(A) \cap \pi_Y^{-1}(B) \in \mathcal{S}
    \]
    since both $\pi_X^{-1}(A), \pi_Y^{-1}(B) \in \mathcal{S}$ and $\mathcal{S}$ is closed under intersections.
    
    Since $\mathcal{M} \otimes \mathcal{N}$ is the $\sigma$-algebra \textit{generated} by rectangles, 
    and $\mathcal{S}$ contains all rectangles and is a $\sigma$-algebra, we have 
    $\mathcal{M} \otimes \mathcal{N} \subseteq \mathcal{S}$.
    
    \bigskip
    \textbf{Step 2: $\mathcal{S} \subseteq \mathcal{M} \otimes \mathcal{N}$.}
    
    For any $A \in \mathcal{M}$:
    \[
        \pi_X^{-1}(A) = A \times Y \in \mathcal{M} \otimes \mathcal{N}
    \]
    since $A \times Y$ is a rectangle (with $B = Y \in \mathcal{N}$).
    
    Similarly, $\pi_Y^{-1}(B) = X \times B \in \mathcal{M} \otimes \mathcal{N}$ for all $B \in \mathcal{N}$.
    
    Thus $\pi_X^{-1}(\mathcal{M}) \cup \pi_Y^{-1}(\mathcal{N}) \subseteq \mathcal{M} \otimes \mathcal{N}$.
    
    Since $\mathcal{S}$ is the \textit{smallest} $\sigma$-algebra containing these sets, 
    and $\mathcal{M} \otimes \mathcal{N}$ is a $\sigma$-algebra containing them, 
    we have $\mathcal{S} \subseteq \mathcal{M} \otimes \mathcal{N}$.
    
    \bigskip
    \textbf{Conclusion:} $\mathcal{M} \otimes \mathcal{N} = \mathcal{S}$.
\end{proof}

\begin{corollary}
    The canonical projections $\pi_X$ and $\pi_Y$ are $(\mathcal{M} \otimes \mathcal{N})$-measurable.
\end{corollary}

\begin{proof}
    By the characterization, $\mathcal{M} \otimes \mathcal{N}$ contains $\pi_X^{-1}(A)$ for all 
    $A \in \mathcal{M}$ and $\pi_Y^{-1}(B)$ for all $B \in \mathcal{N}$.
\end{proof}

\bigskip
\begin{proposition}[Product Borel $\sigma$-Algebra --- General Case]
    For arbitrary topological spaces $X$ and $Y$:
    \[
        \BorelAlg{X} \otimes \BorelAlg{Y} \subseteq \BorelAlg{X \times Y}
    \]
    where $X \times Y$ is given the product topology.
\end{proposition}

\begin{proof}
    The canonical projections $\pi_X: X \times Y \to X$ and $\pi_Y: X \times Y \to Y$ are 
    continuous with respect to the product topology (by definition of product topology).
    
    \bigskip
    \textbf{Step 1: Projections are Borel measurable.}
    
    Since $\pi_X$ is continuous and $\BorelAlg{X}$ is generated by open sets: for any open $U \subseteq X$,
    $\pi_X^{-1}(U)$ is open in $X \times Y$, hence in $\BorelAlg{X \times Y}$.
    
    Thus $\pi_X^{-1}(A) \in \BorelAlg{X \times Y}$ for all $A \in \BorelAlg{X}$.
    Similarly for $\pi_Y$.
    
    \bigskip
    \textbf{Step 2: Rectangles are Borel.}
    
    For any $A \in \BorelAlg{X}$ and $B \in \BorelAlg{Y}$:
    \[
        A \times B = \pi_X^{-1}(A) \cap \pi_Y^{-1}(B) \in \BorelAlg{X \times Y}
    \]
    since $\BorelAlg{X \times Y}$ is closed under intersections.
    
    \bigskip
    \textbf{Step 3: Conclude.}
    
    Since $\BorelAlg{X} \otimes \BorelAlg{Y}$ is the $\sigma$-algebra \textit{generated} by 
    rectangles $\{A \times B\}$, and all rectangles are in $\BorelAlg{X \times Y}$, we have:
    \[
        \BorelAlg{X} \otimes \BorelAlg{Y} \subseteq \BorelAlg{X \times Y}
    \]
\end{proof}

\begin{proposition}[Product Borel $\sigma$-Algebra --- Second Countable Case]
    Let $X$ and $Y$ be second countable topological spaces. Then:
    \[
        \BorelAlg{X \times Y} = \BorelAlg{X} \otimes \BorelAlg{Y}
    \]
\end{proposition}

\begin{proof}
    We already have $\BorelAlg{X} \otimes \BorelAlg{Y} \subseteq \BorelAlg{X \times Y}$ from the 
    previous proposition. It remains to show the reverse inclusion.
    
    \bigskip
    \textbf{Step 1: Identify countable bases.}
    
    Since $X$ is second countable, there exists a countable base $\mathcal{B}_X = \{U_n\}_{n \in \bb{N}}$ 
    for the topology on $X$. Similarly, $Y$ has a countable base $\mathcal{B}_Y = \{V_m\}_{m \in \bb{N}}$.
    
    \bigskip
    \textbf{Step 2: The product topology has a countable base.}
    
    The collection $\mathcal{B} = \{U_n \times V_m \mid n, m \in \bb{N}\}$ is a countable base for the 
    product topology on $X \times Y$. Indeed:
    \begin{itemize}
        \item $\mathcal{B}$ is countable since $\bb{N} \times \bb{N}$ is countable.
        \item For any open set $W \subseteq X \times Y$ and any point $(x, y) \in W$, there exist 
              open sets $U \subseteq X$ and $V \subseteq Y$ with $(x, y) \in U \times V \subseteq W$ 
              (by definition of product topology). Since $\mathcal{B}_X$ and $\mathcal{B}_Y$ are bases, 
              there exist $U_n \ni x$ with $U_n \subseteq U$ and $V_m \ni y$ with $V_m \subseteq V$. 
              Thus $(x, y) \in U_n \times V_m \subseteq W$.
    \end{itemize}
    
    \bigskip
    \textbf{Step 3: Open sets are countable unions of base elements.}
    
    Every open set $W \subseteq X \times Y$ can be written as:
    \[
        W = \bigcup_{(n, m) \in I} (U_n \times V_m)
    \]
    for some countable index set $I \subseteq \bb{N} \times \bb{N}$.
    
    \bigskip
    \textbf{Step 4: Base elements are in $\BorelAlg{X} \otimes \BorelAlg{Y}$.}
    
    Each $U_n$ is open in $X$, hence $U_n \in \BorelAlg{X}$. Similarly, $V_m \in \BorelAlg{Y}$. 
    Therefore:
    \[
        U_n \times V_m \in \BorelAlg{X} \otimes \BorelAlg{Y}
    \]
    since this is precisely a measurable rectangle.
    
    \bigskip
    \textbf{Step 5: Open sets are in $\BorelAlg{X} \otimes \BorelAlg{Y}$.}
    
    Since every open set $W$ is a countable union of sets in $\BorelAlg{X} \otimes \BorelAlg{Y}$, 
    and $\sigma$-algebras are closed under countable unions, we have:
    \[
        W \in \BorelAlg{X} \otimes \BorelAlg{Y}
    \]
    for every open $W \subseteq X \times Y$.
    
    \bigskip
    \textbf{Step 6: Conclude.}
    
    The Borel $\sigma$-algebra $\BorelAlg{X \times Y}$ is generated by open sets in $X \times Y$. 
    Since all open sets are in $\BorelAlg{X} \otimes \BorelAlg{Y}$ (which is a $\sigma$-algebra), 
    we have:
    \[
        \BorelAlg{X \times Y} = \GenSigmaAlg{\text{open sets in } X \times Y} \subseteq \BorelAlg{X} \otimes \BorelAlg{Y}
    \]
    
    Combined with the previous proposition: $\BorelAlg{X \times Y} = \BorelAlg{X} \otimes \BorelAlg{Y}$.
\end{proof}

\begin{remark}
    The reverse inclusion $\BorelAlg{X \times Y} \subseteq \BorelAlg{X} \otimes \BorelAlg{Y}$ may fail 
    for general topological spaces that are not second countable. The key to the proof above is that 
    open sets in the product are \emph{countable} unions of rectangles---this uses second countability 
    essentially.
\end{remark}

\bigskip
\begin{proposition}
    $\BorelAlg{\bb{R}^{n+m}} = \BorelAlg{\bb{R}^n} \otimes \BorelAlg{\bb{R}^m}$.
\end{proposition}

\begin{proof}
    Let $\mathcal{R} = \{A \times B \mid A \in \BorelAlg{\bb{R}^n}, B \in \BorelAlg{\bb{R}^m}\}$. 
    We show $\GenSigmaAlg{\mathcal{R}} = \BorelAlg{\bb{R}^{n+m}}$.

    \bigskip
    ($\subseteq$) The projections $\pi_1: \bb{R}^{n+m} \to \bb{R}^n$ and 
    $\pi_2: \bb{R}^{n+m} \to \bb{R}^m$ are continuous, hence $(\BorelAlg{\bb{R}^{n+m}})$-measurable. 
    For $A \in \BorelAlg{\bb{R}^n}$, $B \in \BorelAlg{\bb{R}^m}$: 
    $A \times B = \pi_1^{-1}(A) \cap \pi_2^{-1}(B) \in \BorelAlg{\bb{R}^{n+m}}$.

    \bigskip
    ($\supseteq$) Open sets in $\bb{R}^{n+m}$ are countable unions of open rectangles. 
    Each open rectangle is in $\GenSigmaAlg{\mathcal{R}}$ since open sets are Borel.
\end{proof}

\section{Sections of Sets and Functions}

\begin{intuition}
    To construct the product measure, we need to ``slice'' sets in $X \times Y$ along one 
    coordinate. For a set $E \subseteq X \times Y$, we can fix an $x \in X$ and ask: 
    ``Which $y$ values are paired with $x$ in $E$?'' This gives the $x$-section $E_x \subseteq Y$.
    Similarly, fixing $y$ gives the $y$-section $E^y \subseteq X$.
    
    The key insight is that if $E$ is measurable in the product $\sigma$-algebra, then 
    its sections are measurable in the respective factor $\sigma$-algebras. This allows us to 
    integrate over sections to construct the product measure.
\end{intuition}

\bigskip
\begin{definition}[Sections of Sets]
    Let $E \subseteq X \times Y$. For $x \in X$ and $y \in Y$, define:
    \begin{itemize}
        \item The \textbf{$x$-section} (or \textbf{horizontal section}) of $E$:
        \[
            E_x = \{y \in Y \mid (x, y) \in E\} \subseteq Y
        \]
        \item The \textbf{$y$-section} (or \textbf{vertical section}) of $E$:
        \[
            E^y = \{x \in X \mid (x, y) \in E\} \subseteq X
        \]
    \end{itemize}
\end{definition}

\begin{example}
    For a rectangle $A \times B$ with $A \subseteq X$ and $B \subseteq Y$:
    \[
        (A \times B)_x = \begin{cases} B & \text{if } x \in A \\ \emptyset & \text{if } x \notin A \end{cases}
        \qquad
        (A \times B)^y = \begin{cases} A & \text{if } y \in B \\ \emptyset & \text{if } y \notin B \end{cases}
    \]
\end{example}

\bigskip
\begin{definition}[Sections of Functions]
    Let $f: X \times Y \to \overline{\bb{R}}$. For $x \in X$ and $y \in Y$, define:
    \begin{itemize}
        \item The \textbf{$x$-section} of $f$: $f_x: Y \to \overline{\bb{R}}$, $f_x(y) = f(x, y)$
        \item The \textbf{$y$-section} of $f$: $f^y: X \to \overline{\bb{R}}$, $f^y(x) = f(x, y)$
    \end{itemize}
\end{definition}

\begin{remark}
    For a set $E \subseteq X \times Y$ and its characteristic function $\CharFunc{E}$:
    \[
        (\CharFunc{E})_x = \CharFunc{E_x} \qquad \text{and} \qquad (\CharFunc{E})^y = \CharFunc{E^y}
    \]
\end{remark}

\bigskip
\begin{proposition}[Properties of Sections]
    Let $E, F \subseteq X \times Y$ and $(E_n)_{n \in \bb{N}}$ be a sequence of subsets. Then:
    \begin{enumerate}
        \item $(E^c)_x = (E_x)^c$ and $(E^c)^y = (E^y)^c$
        \item $(E \cup F)_x = E_x \cup F_x$ and $(E \cup F)^y = E^y \cup F^y$
        \item $(E \cap F)_x = E_x \cap F_x$ and $(E \cap F)^y = E^y \cap F^y$
        \item $\left(\bigcup_{n \in \bb{N}} E_n\right)_x = \bigcup_{n \in \bb{N}} (E_n)_x$ and similarly for $y$-sections
        \item $\left(\bigcap_{n \in \bb{N}} E_n\right)_x = \bigcap_{n \in \bb{N}} (E_n)_x$ and similarly for $y$-sections
    \end{enumerate}
\end{proposition}

\begin{proof}
    These follow directly from the definitions. We prove (1) and (4), the others being similar.
    
    \bigskip
    \textbf{Property 1:}
    \begin{align*}
        (E^c)_x &= \{y \in Y \mid (x, y) \in E^c\} \\
                &= \{y \in Y \mid (x, y) \notin E\} \\
                &= \{y \in Y \mid y \notin E_x\} \\
                &= Y \setminus E_x = (E_x)^c
    \end{align*}
    
    \bigskip
    \textbf{Property 4:}
    \begin{align*}
        \left(\bigcup_{n \in \bb{N}} E_n\right)_x &= \left\{y \in Y \mid (x, y) \in \bigcup_{n \in \bb{N}} E_n\right\} \\
        &= \{y \in Y \mid \Exists :{n \in \bb{N}}{(x, y) \in E_n}\} \\
        &= \{y \in Y \mid \Exists :{n \in \bb{N}}{y \in (E_n)_x}\} \\
        &= \bigcup_{n \in \bb{N}} (E_n)_x
    \end{align*}
\end{proof}

\bigskip
\begin{theorem}[Measurability of Sections]
    Let $(X, \mathcal{M})$ and $(Y, \mathcal{N})$ be measurable spaces. If $E \in \mathcal{M} \otimes \mathcal{N}$, then:
    \begin{enumerate}
        \item $\Forall :{x \in X}{E_x \in \mathcal{N}}$ (all $x$-sections are $\mathcal{N}$-measurable)
        \item $\Forall :{y \in Y}{E^y \in \mathcal{M}}$ (all $y$-sections are $\mathcal{M}$-measurable)
    \end{enumerate}
\end{theorem}

\begin{proof}
    We prove (1); the proof of (2) is symmetric.
    
    \bigskip
    \textbf{Strategy:} We use the $\sigma$-algebra induction principle. Define:
    \[
        \mathcal{G} = \{E \in \mathcal{M} \otimes \mathcal{N} \mid \Forall :{x \in X}{E_x \in \mathcal{N}}\}
    \]
    We show $\mathcal{G}$ is a $\sigma$-algebra containing all rectangles, hence $\mathcal{G} = \mathcal{M} \otimes \mathcal{N}$.
    
    \bigskip
    \textbf{Step 1: Rectangles are in $\mathcal{G}$.}
    
    Let $E = A \times B$ with $A \in \mathcal{M}$, $B \in \mathcal{N}$. Then:
    \[
        E_x = (A \times B)_x = \begin{cases} B & \text{if } x \in A \\ \emptyset & \text{if } x \notin A \end{cases}
    \]
    Since $B \in \mathcal{N}$ and $\emptyset \in \mathcal{N}$, we have $E_x \in \mathcal{N}$ for all $x$.
    Thus $A \times B \in \mathcal{G}$.
    
    \bigskip
    \textbf{Step 2: $\mathcal{G}$ is a $\sigma$-algebra.}
    
    \textit{Contains $\emptyset$:} $\emptyset_x = \emptyset \in \mathcal{N}$ for all $x$, so $\emptyset \in \mathcal{G}$.
    
    \textit{Closed under complements:} Let $E \in \mathcal{G}$. Then $(E^c)_x = (E_x)^c \in \mathcal{N}$ 
    (since $\mathcal{N}$ is closed under complements). Thus $E^c \in \mathcal{G}$.
    
    \textit{Closed under countable unions:} Let $(E_n)_{n \in \bb{N}} \subseteq \mathcal{G}$. Then:
    \[
        \left(\bigcup_{n \in \bb{N}} E_n\right)_x = \bigcup_{n \in \bb{N}} (E_n)_x \in \mathcal{N}
    \]
    since each $(E_n)_x \in \mathcal{N}$ and $\mathcal{N}$ is closed under countable unions.
    
    \bigskip
    \textbf{Step 3: Conclude.}
    
    Since $\mathcal{G}$ is a $\sigma$-algebra containing all rectangles $\{A \times B\}$, and 
    $\mathcal{M} \otimes \mathcal{N}$ is the smallest such $\sigma$-algebra:
    \[
        \mathcal{M} \otimes \mathcal{N} \subseteq \mathcal{G}
    \]
    But $\mathcal{G} \subseteq \mathcal{M} \otimes \mathcal{N}$ by definition. Hence $\mathcal{G} = \mathcal{M} \otimes \mathcal{N}$.
\end{proof}

\bigskip
\section{Product Measures}

\begin{intuition}
    Given measure spaces $(X, \mathcal{M}, \mu)$ and $(Y, \mathcal{N}, \nu)$, we want to define a 
    measure on the product $\sigma$-algebra that assigns measure $\mu(A) \cdot \nu(B)$ to each 
    rectangle $A \times B$. 
    
    The construction uses sections: for a measurable set $E \subseteq X \times Y$, we integrate the 
    ``slice measures'' $\nu(E_x)$ over $x$. The key technical challenges are:
    \begin{enumerate}
        \item Showing $x \mapsto \nu(E_x)$ is measurable
        \item Showing this construction yields a measure
        \item Proving uniqueness
    \end{enumerate}
    
    We use $\sigma$-finiteness to ensure uniqueness via the $\pi$-$\lambda$ theorem.
\end{intuition}

\bigskip
\begin{lemma}[Measurability of Section Measures]
    Let $(X, \mathcal{M}, \mu)$ and $(Y, \mathcal{N}, \nu)$ be $\sigma$-finite measure spaces. 
    For $E \in \mathcal{M} \otimes \mathcal{N}$:
    \begin{enumerate}
        \item The function $\varphi: X \to [0, \infty]$, $\varphi(x) = \nu(E_x)$ is $\mathcal{M}$-measurable
        \item The function $\psi: Y \to [0, \infty]$, $\psi(y) = \mu(E^y)$ is $\mathcal{N}$-measurable
    \end{enumerate}
\end{lemma}

\begin{proof}
    We prove (1); the proof of (2) is symmetric. We proceed in two stages: first for finite 
    measures, then extend to $\sigma$-finite.
    
    \textbf{Case 1: $\nu$ is a finite measure.}
    
    Define:
    \[
        \mathcal{D} = \{E \in \mathcal{M} \otimes \mathcal{N} \mid x \mapsto \nu(E_x) \text{ is } \mathcal{M}\text{-measurable}\}
    \]
    
    We show $\mathcal{D}$ is a $\lambda$-system containing all rectangles.
    
    \bigskip
    \textit{$\mathcal{D}$ contains rectangles:}
    
    For $E = A \times B$ with $A \in \mathcal{M}$, $B \in \mathcal{N}$:
    \[
        \nu(E_x) = \nu((A \times B)_x) = \begin{cases} \nu(B) & \text{if } x \in A \\ 0 & \text{if } x \notin A \end{cases}
        = \nu(B) \cdot \CharFunc{A}(x)
    \]
    This is $\mathcal{M}$-measurable since $\CharFunc{A}$ is measurable.
    
    \bigskip
    \textit{$\mathcal{D}$ is a $\lambda$-system:}
    
    \textbf{Axiom 1:} $(X \times Y)_x = Y$, so $\nu((X \times Y)_x) = \nu(Y) < \infty$ (constant function). 
    Thus $X \times Y \in \mathcal{D}$.
    
    \textbf{Axiom 2:} Let $E, F \in \mathcal{D}$ with $E \subseteq F$. Then:
    \[
        \nu((F \setminus E)_x) = \nu(F_x \setminus E_x) = \nu(F_x) - \nu(E_x)
    \]
    (using that $E_x \subseteq F_x$ and $\nu$ is finite). The right-hand side is the difference of two 
    measurable functions, hence measurable. Thus $F \setminus E \in \mathcal{D}$.
    
    \textbf{Axiom 3:} Let $(E_n)_{n \in \bb{N}} \subseteq \mathcal{D}$ with $E_n \nearrow E$. 
    Then $(E_n)_x \nearrow E_x$ for each $x$. By continuity from below:
    \[
        \nu(E_x) = \lim_{n \to \infty} \nu((E_n)_x)
    \]
    The right-hand side is the pointwise limit of measurable functions, hence measurable. 
    Thus $E \in \mathcal{D}$.
    
    \bigskip
    \textit{Apply Dynkin $\pi$-$\lambda$ theorem:}
    
    The rectangles $\mathcal{R} = \{A \times B \mid A \in \mathcal{M}, B \in \mathcal{N}\}$ form a $\pi$-system 
    (since $(A_1 \times B_1) \cap (A_2 \times B_2) = (A_1 \cap A_2) \times (B_1 \cap B_2)$).
    
    We have shown $\mathcal{R} \subseteq \mathcal{D}$ and $\mathcal{D}$ is a $\lambda$-system.
    By the $\pi$-$\lambda$ theorem:
    \[
        \mathcal{M} \otimes \mathcal{N} = \GenSigmaAlg{\mathcal{R}} = \GenLambda{\mathcal{R}} \subseteq \mathcal{D}
    \]
    Hence $\mathcal{D} = \mathcal{M} \otimes \mathcal{N}$.
    
    \bigskip
    \textbf{Case 2: $\nu$ is $\sigma$-finite.}
    
    Write $Y = \bigcup_{n \in \bb{N}} Y_n$ with $Y_n \nearrow$ and $\nu(Y_n) < \infty$.
    
    For each $n$, the restricted measure $\nu_n = \nu|_{Y_n}$ is finite. Define $E_n = E \cap (X \times Y_n)$.
    By Case 1, $x \mapsto \nu((E_n)_x)$ is measurable.
    
    Note that $(E_n)_x = E_x \cap Y_n \nearrow E_x$ as $n \to \infty$. By continuity from below:
    \[
        \nu(E_x) = \lim_{n \to \infty} \nu(E_x \cap Y_n) = \lim_{n \to \infty} \nu((E_n)_x)
    \]
    This is the limit of measurable functions, hence measurable.
\end{proof}

\bigskip
\begin{theorem}[Construction of Product Measure via $x$-Sections]
    Let $(X, \mathcal{M}, \mu)$ and $(Y, \mathcal{N}, \nu)$ be $\sigma$-finite measure spaces. 
    Define $\rho_x: \mathcal{M} \otimes \mathcal{N} \to [0, \infty]$ by:
    \[
        \rho_x(E) = \Integral_X \nu(E_x) \, d\mu(x)
    \]
    Then $\rho_x$ is a measure on $\mathcal{M} \otimes \mathcal{N}$ satisfying:
    \[
        \rho_x(A \times B) = \mu(A) \cdot \nu(B)
    \]
    for all $A \in \mathcal{M}$, $B \in \mathcal{N}$.
\end{theorem}

\begin{proof}
    We verify the measure axioms for $\rho_x$.
    
    \bigskip
    \textbf{Step 1: $\rho_x(\emptyset) = 0$.}
    
    We have $\emptyset_x = \emptyset$ for all $x$, so $\nu(\emptyset_x) = 0$. Thus:
    \[
        \rho_x(\emptyset) = \Integral_X 0 \, d\mu(x) = 0
    \]
    
    \bigskip
    \textbf{Step 2: $\sigma$-additivity.}
    
    Let $(E_n)_{n \in \bb{N}}$ be pairwise disjoint sets in $\mathcal{M} \otimes \mathcal{N}$. 
    Let $E = \bigsqcup_{n \in \bb{N}} E_n$.
    
    For each $x \in X$, the sections $(E_n)_x$ are pairwise disjoint in $Y$ (since if $y \in (E_i)_x \cap (E_j)_x$ 
    with $i \neq j$, then $(x, y) \in E_i \cap E_j = \emptyset$, contradiction).
    
    Moreover, $E_x = \bigsqcup_{n \in \bb{N}} (E_n)_x$. By $\sigma$-additivity of $\nu$:
    \[
        \nu(E_x) = \sum_{n \in \bb{N}} \nu((E_n)_x)
    \]
    
    Now by the Monotone Convergence Theorem (applied to the partial sums):
    \begin{align*}
        \rho_x(E) &= \Integral_X \nu(E_x) \, d\mu(x) \\
        &= \Integral_X \sum_{n \in \bb{N}} \nu((E_n)_x) \, d\mu(x) \\
        &= \sum_{n \in \bb{N}} \Integral_X \nu((E_n)_x) \, d\mu(x) \quad \text{(by MCT)} \\
        &= \sum_{n \in \bb{N}} \rho_x(E_n)
    \end{align*}
    
    \bigskip
    \textbf{Step 3: Value on rectangles.}
    
    For $E = A \times B$ with $A \in \mathcal{M}$, $B \in \mathcal{N}$:
    \[
        \nu(E_x) = \nu((A \times B)_x) = \nu(B) \cdot \CharFunc{A}(x)
    \]
    Thus:
    \[
        \rho_x(A \times B) = \Integral_X \nu(B) \cdot \CharFunc{A}(x) \, d\mu(x) = \nu(B) \cdot \mu(A)
    \]
\end{proof}

\bigskip
\begin{theorem}[Construction of Product Measure via $y$-Sections]
    Let $(X, \mathcal{M}, \mu)$ and $(Y, \mathcal{N}, \nu)$ be $\sigma$-finite measure spaces. 
    Define $\rho_y: \mathcal{M} \otimes \mathcal{N} \to [0, \infty]$ by:
    \[
        \rho_y(E) = \Integral_Y \mu(E^y) \, d\nu(y)
    \]
    Then $\rho_y$ is a measure on $\mathcal{M} \otimes \mathcal{N}$ satisfying:
    \[
        \rho_y(A \times B) = \mu(A) \cdot \nu(B)
    \]
\end{theorem}

\begin{proof}
    Symmetric to the proof for $\rho_x$.
\end{proof}

\bigskip
\begin{theorem}[Uniqueness of Product Measure]
    Let $(X, \mathcal{M}, \mu)$ and $(Y, \mathcal{N}, \nu)$ be $\sigma$-finite measure spaces. 
    Then $\rho_x = \rho_y$. Moreover, there exists a \textbf{unique} measure $\mu \otimes \nu$ on 
    $\mathcal{M} \otimes \mathcal{N}$ such that:
    \[
        (\mu \otimes \nu)(A \times B) = \mu(A) \cdot \nu(B)
    \]
    for all $A \in \mathcal{M}$, $B \in \mathcal{N}$.
\end{theorem}

\begin{proof}
    The rectangles $\mathcal{R} = \{A \times B \mid A \in \mathcal{M}, B \in \mathcal{N}\}$ form a 
    $\pi$-system generating $\mathcal{M} \otimes \mathcal{N}$:
    \[
        (A_1 \times B_1) \cap (A_2 \times B_2) = (A_1 \cap A_2) \times (B_1 \cap B_2) \in \mathcal{R}
    \]
    
    Both $\rho_x$ and $\rho_y$ are $\sigma$-finite measures on $\mathcal{M} \otimes \mathcal{N}$ 
    (since $(X, \mu)$ and $(Y, \nu)$ are $\sigma$-finite) that agree on $\mathcal{R}$:
    \[
        \rho_x(A \times B) = \mu(A) \cdot \nu(B) = \rho_y(A \times B)
    \]
    
    By the Uniqueness Theorem for $\sigma$-finite measures agreeing on a $\pi$-system:
    $\rho_x = \rho_y$ on all of $\mathcal{M} \otimes \mathcal{N}$.
    
    The same argument shows any measure satisfying the rectangle property must equal $\rho_x = \rho_y$.
\end{proof}

\bigskip
\begin{definition}[Product Measure]
    The \textbf{product measure} $\mu \otimes \nu$ on $(X \times Y, \mathcal{M} \otimes \mathcal{N})$ 
    is the unique measure satisfying:
    \[
        (\mu \otimes \nu)(A \times B) = \mu(A) \cdot \nu(B)
    \]
    Equivalently:
    \[
        (\mu \otimes \nu)(E) = \Integral_X \nu(E_x) \, d\mu(x) = \Integral_Y \mu(E^y) \, d\nu(y)
    \]
\end{definition}

\section{Tonelli's Theorem}

\begin{intuition}
    Tonelli's Theorem allows us to evaluate integrals over product spaces by iterating 
    one-dimensional integrals. The key insight is that for \textbf{non-negative} measurable 
    functions, no integrability assumptions are needed---the iterated integrals always 
    equal the product integral (though all three may be infinite).
    
    The theorem proceeds by:
    \begin{enumerate}
        \item Verifying for characteristic functions of measurable sets
        \item Extending to simple functions by linearity
        \item Extending to general non-negative functions via monotone approximation
    \end{enumerate}
    
    The non-negativity is essential: it allows us to apply MCT freely without worrying 
    about $\infty - \infty$ issues.
\end{intuition}

\bigskip
\begin{theorem}[Tonelli]
    Let $(X, \mathcal{M}, \mu)$ and $(Y, \mathcal{N}, \nu)$ be $\sigma$-finite measure spaces.
    Let $f: X \times Y \to [0, \infty]$ be $(\mathcal{M} \otimes \mathcal{N})$-measurable. Then:
    \begin{enumerate}
        \item For each $x \in X$, the $x$-section $f_x: Y \to [0, \infty]$ is $\mathcal{N}$-measurable
        \item For each $y \in Y$, the $y$-section $f^y: X \to [0, \infty]$ is $\mathcal{M}$-measurable
        \item The function $x \mapsto \Integral_Y f(x, y) \, d\nu(y)$ is $\mathcal{M}$-measurable
        \item The function $y \mapsto \Integral_X f(x, y) \, d\mu(x)$ is $\mathcal{N}$-measurable
        \item The iterated integrals equal the product integral:
        \[
            \Integral_{X \times Y} f \, d(\mu \otimes \nu) = \Integral_X \left(\Integral_Y f(x, y) \, d\nu(y)\right) d\mu(x) = \Integral_Y \left(\Integral_X f(x, y) \, d\mu(x)\right) d\nu(y)
        \]
    \end{enumerate}
\end{theorem}

\begin{proof}
    We prove the theorem in four steps, building from simple cases to the general result.
    
    \bigskip
    \textbf{Step 1: Measurability of sections.}
    
    \textit{Claim:} For any $(\mathcal{M} \otimes \mathcal{N})$-measurable function $f: X \times Y \to [0, \infty]$, 
    each section $f_x$ is $\mathcal{N}$-measurable and each section $f^y$ is $\mathcal{M}$-measurable.
    
    \textit{Proof of claim:} For any $a \in [0, \infty]$:
    \[
        f_x^{-1}([0, a]) = \{y \in Y \mid f(x, y) \leq a\} = \{y \in Y \mid (x, y) \in f^{-1}([0, a])\}
    \]
    Since $f$ is $(\mathcal{M} \otimes \mathcal{N})$-measurable, $f^{-1}([0, a]) \in \mathcal{M} \otimes \mathcal{N}$.
    By the Measurability of Sections theorem, $(f^{-1}([0, a]))_x \in \mathcal{N}$.
    But $(f^{-1}([0, a]))_x = f_x^{-1}([0, a])$, so $f_x$ is $\mathcal{N}$-measurable.
    
    The argument for $f^y$ is symmetric.
    
    \bigskip
    \textbf{Step 2: Characteristic functions of measurable sets.}
    
    Let $E \in \mathcal{M} \otimes \mathcal{N}$ and $f = \CharFunc{E}$.
    
    \textit{Computing the inner integral:}
    For each $x \in X$:
    \[
        \Integral_Y \CharFunc{E}(x, y) \, d\nu(y) = \Integral_Y \CharFunc{E_x}(y) \, d\nu(y) = \nu(E_x)
    \]
    where $E_x = \{y \in Y \mid (x, y) \in E\}$ is the $x$-section of $E$.
    
    \textit{Measurability of the inner integral:}
    By the Measurability of Section Measures lemma, $x \mapsto \nu(E_x)$ is $\mathcal{M}$-measurable.
    
    \textit{Computing the iterated integral:}
    \[
        \Integral_X \left(\Integral_Y \CharFunc{E}(x, y) \, d\nu(y)\right) d\mu(x) = \Integral_X \nu(E_x) \, d\mu(x) = (\mu \otimes \nu)(E)
    \]
    by definition of the product measure.
    
    \textit{Product integral:}
    \[
        \Integral_{X \times Y} \CharFunc{E} \, d(\mu \otimes \nu) = (\mu \otimes \nu)(E)
    \]
    
    Thus the iterated integral equals the product integral for characteristic functions.
    
    By symmetry, the same holds integrating first over $X$:
    \[
        \Integral_Y \left(\Integral_X \CharFunc{E}(x, y) \, d\mu(x)\right) d\nu(y) = \Integral_Y \mu(E^y) \, d\nu(y) = (\mu \otimes \nu)(E)
    \]
    
    \bigskip
    \textbf{Step 3: Simple functions.}
    
    Let $s = \sum_{i=1}^n c_i \CharFunc{E_i}$ be a simple function in canonical form, where 
    $c_i \geq 0$ and $E_i \in \mathcal{M} \otimes \mathcal{N}$ are pairwise disjoint.
    
    \textit{Measurability:} By Step 1, each section $s_x$ is $\mathcal{N}$-measurable.
    
    \textit{Inner integral:} By linearity of the integral:
    \[
        \Integral_Y s(x, y) \, d\nu(y) = \sum_{i=1}^n c_i \Integral_Y \CharFunc{E_i}(x, y) \, d\nu(y) = \sum_{i=1}^n c_i \nu((E_i)_x)
    \]
    This is a finite sum of $\mathcal{M}$-measurable functions (by Step 2), hence $\mathcal{M}$-measurable.
    
    \textit{Iterated integral:} By linearity of the integral (justified since $c_i \geq 0$):
    \begin{align*}
        \Integral_X \left(\Integral_Y s \, d\nu\right) d\mu 
        &= \Integral_X \sum_{i=1}^n c_i \nu((E_i)_x) \, d\mu(x) \\
        &= \sum_{i=1}^n c_i \Integral_X \nu((E_i)_x) \, d\mu(x) \\
        &= \sum_{i=1}^n c_i (\mu \otimes \nu)(E_i) \quad \text{(by Step 2)} \\
        &= \Integral_{X \times Y} s \, d(\mu \otimes \nu)
    \end{align*}
    
    By symmetry, the same holds integrating first over $X$.
    
    \bigskip
    \textbf{Step 4: General non-negative measurable functions.}
    
    Let $f: X \times Y \to [0, \infty]$ be $(\mathcal{M} \otimes \mathcal{N})$-measurable.
    
    By the Simple Function Approximation Theorem, there exists a sequence $(s_n)_{n \in \bb{N}}$ 
    of simple functions with $0 \leq s_n \nearrow f$ pointwise.
    
    \textit{Measurability of the inner integral:}
    For each $x$, we have $0 \leq (s_n)_x \nearrow f_x$ pointwise on $Y$. By MCT:
    \[
        \Integral_Y f_x \, d\nu = \lim_{n \to \infty} \Integral_Y (s_n)_x \, d\nu
    \]
    Define $\varphi_n(x) = \Integral_Y (s_n)_x \, d\nu$ and $\varphi(x) = \Integral_Y f_x \, d\nu$.
    
    Each $\varphi_n$ is $\mathcal{M}$-measurable (by Step 3).
    Since $s_n \nearrow f$, we have $\varphi_n(x) \nearrow \varphi(x)$ for each $x$ (by MCT).
    Thus $\varphi = \lim_n \varphi_n$ is $\mathcal{M}$-measurable (as the pointwise limit of measurable functions).
    
    \textit{Equality of integrals:}
    By Step 3, for each $n$:
    \[
        \Integral_{X \times Y} s_n \, d(\mu \otimes \nu) = \Integral_X \varphi_n \, d\mu
    \]
    
    Taking limits as $n \to \infty$:
    \begin{itemize}
        \item \textbf{Left side:} By MCT on $(X \times Y, \mu \otimes \nu)$:
        \[
            \lim_{n \to \infty} \Integral_{X \times Y} s_n \, d(\mu \otimes \nu) = \Integral_{X \times Y} f \, d(\mu \otimes \nu)
        \]
        
        \item \textbf{Right side:} Since $\varphi_n \nearrow \varphi$, by MCT on $(X, \mu)$:
        \[
            \lim_{n \to \infty} \Integral_X \varphi_n \, d\mu = \Integral_X \varphi \, d\mu = \Integral_X \left(\Integral_Y f \, d\nu\right) d\mu
        \]
    \end{itemize}
    
    Therefore:
    \[
        \Integral_{X \times Y} f \, d(\mu \otimes \nu) = \Integral_X \left(\Integral_Y f(x, y) \, d\nu(y)\right) d\mu(x)
    \]
    
    By symmetry (exchanging roles of $X$ and $Y$):
    \[
        \Integral_{X \times Y} f \, d(\mu \otimes \nu) = \Integral_Y \left(\Integral_X f(x, y) \, d\mu(x)\right) d\nu(y)
    \]
    
    This completes the proof.
\end{proof}

\begin{corollary}
    If one of the iterated integrals of a non-negative measurable function is finite, 
    then all three integrals (two iterated, one product) are equal and finite.
\end{corollary}

\begin{remark}
    Tonelli's Theorem does \textbf{not} require integrability. The equality holds even when 
    all integrals are infinite. However, when applying the theorem to functions that may 
    change sign, we must first apply it to $|f|$ to verify integrability before using Fubini.
\end{remark}

\bigskip
\section{Fubini's Theorem}

\begin{intuition}
    Fubini's Theorem extends Tonelli to \textbf{integrable} functions that may take both 
    positive and negative values. The key difference is that we need the integrability 
    condition $f \in \Lp{1}(\mu \otimes \nu)$ to ensure the iterated integrals are well-defined 
    (avoiding $\infty - \infty$).
    
    The strategy is:
    \begin{enumerate}
        \item Apply Tonelli to $|f|$ to establish that sections are integrable a.e.
        \item Decompose $f = f^+ - f^-$ and apply Tonelli to each part
        \item Subtract, using that both parts have finite integrals
    \end{enumerate}
\end{intuition}

\bigskip
\begin{theorem}[Fubini]
    Let $(X, \mathcal{M}, \mu)$ and $(Y, \mathcal{N}, \nu)$ be $\sigma$-finite measure spaces.
    Let $f: X \times Y \to \overline{\bb{R}}$ be $(\mathcal{M} \otimes \mathcal{N})$-measurable with 
    $f \in \Lp{1}(\mu \otimes \nu)$. Then:
    \begin{enumerate}
        \item There exists a $\mu$-null set $N \subseteq X$ such that for all $x \in X \setminus N$: 
              the section $f_x \in \Lp{1}(\nu)$
        \item There exists a $\nu$-null set $M \subseteq Y$ such that for all $y \in Y \setminus M$: 
              the section $f^y \in \Lp{1}(\mu)$
        \item The function 
              \[
                  h(x) = \begin{cases}
                      \Integral_Y f(x, y) \, d\nu(y) & \text{if } x \in X \setminus N \\
                      0 & \text{if } x \in N
                  \end{cases}
              \]
              is in $\Lp{1}(\mu)$ (and the choice of value on $N$ does not affect the integral)
        \item Similarly for $y \mapsto \Integral_X f(x, y) \, d\mu(x)$
        \item The iterated integrals equal the product integral:
        \[
            \Integral_{X \times Y} f \, d(\mu \otimes \nu) = \Integral_{X \setminus N} \left(\Integral_Y f(x, y) \, d\nu(y)\right) d\mu(x) = \Integral_{Y \setminus M} \left(\Integral_X f(x, y) \, d\mu(x)\right) d\nu(y)
        \]
    \end{enumerate}
\end{theorem}

\begin{remark}[Completeness Issue]
    The product of complete measure spaces need not be complete. Even if $(X, \mathcal{M}, \mu)$ 
    and $(Y, \mathcal{N}, \nu)$ are complete, the product measure space $(X \times Y, \mathcal{M} \otimes \mathcal{N}, \mu \otimes \nu)$ 
    may fail to be complete. 
    
    This matters because the section $f_x$ may fail to be $\mathcal{N}$-measurable for $x$ in a 
    $\mu$-null set $N$---even though $f$ is $(\mathcal{M} \otimes \mathcal{N})$-measurable. 
    The issue is that if $N \times Y$ is a null set in the product, its ``slices'' $\{x\} \times Y$ 
    for $x \in N$ may contain non-measurable subsets of $Y$ (since the product $\sigma$-algebra 
    does not include all subsets of null sets unless completed).
    
    This is why we must state the theorem with explicit null sets $N$ and $M$, and integrate 
    over $X \setminus N$ and $Y \setminus M$ respectively.
\end{remark}

\begin{proof}
    We proceed in four steps, carefully handling the exceptional null sets.
    
    \bigskip
    \textbf{Step 1: Apply Tonelli to $|f|$.}
    
    Since $|f|: X \times Y \to [0, \infty]$ is non-negative and $(\mathcal{M} \otimes \mathcal{N})$-measurable, 
    Tonelli's Theorem applies:
    \[
        \Integral_{X \times Y} |f| \, d(\mu \otimes \nu) = \Integral_X \left(\Integral_Y |f(x, y)| \, d\nu(y)\right) d\mu(x)
    \]
    
    By hypothesis, $f \in \Lp{1}(\mu \otimes \nu)$, so:
    \[
        \Integral_{X \times Y} |f| \, d(\mu \otimes \nu) < \infty
    \]
    
    \bigskip
    \textbf{Step 2: Identify the exceptional null set and establish a.e.\ integrability.}
    
    Define $g: X \to [0, \infty]$ by:
    \[
        g(x) = \Integral_Y |f(x, y)| \, d\nu(y) = \Integral_Y |f_x| \, d\nu
    \]
    
    By Tonelli, $g$ is $\mathcal{M}$-measurable and:
    \[
        \Integral_X g \, d\mu = \Integral_{X \times Y} |f| \, d(\mu \otimes \nu) < \infty
    \]
    
    Define the exceptional null set:
    \[
        N = \{x \in X \mid g(x) = \infty\} = \left\{x \in X \mid \Integral_Y |f_x| \, d\nu = \infty\right\}
    \]
    
    \textit{Claim:} $\mu(N) = 0$.
    
    \textit{Proof:} If $\mu(N) > 0$, then:
    \[
        \Integral_X g \, d\mu \geq \Integral_N g \, d\mu = \Integral_N \infty \, d\mu = \infty \cdot \mu(N) = \infty
    \]
    contradicting $\Integral_X g \, d\mu < \infty$.
    
    Therefore, for all $x \in X \setminus N$:
    \[
        \Integral_Y |f(x, y)| \, d\nu(y) < \infty
    \]
    which means $f_x \in \Lp{1}(\nu)$ for all $x \in X \setminus N$.
    
    \textit{Important:} We claim integrability for \textbf{all} $x \in X \setminus N$, not merely 
    ``$\mu$-a.e.\ $x$''. The distinction matters because the inner integral is well-defined 
    precisely on $X \setminus N$.
    
    By symmetry, there exists a $\nu$-null set $M \subseteq Y$ such that $f^y \in \Lp{1}(\mu)$ 
    for all $y \in Y \setminus M$.
    
    \bigskip
    \textbf{Step 3: Define the inner integral function on $X \setminus N$.}
    
    For $x \in X \setminus N$, define:
    \[
        h(x) = \Integral_Y f(x, y) \, d\nu(y)
    \]
    This is well-defined since $f_x \in \Lp{1}(\nu)$.
    
    Extend $h$ to all of $X$ by setting $h(x) = 0$ for $x \in N$. (The specific value on $N$ 
    is irrelevant since $\mu(N) = 0$.)
    
    \textit{$h$ is $\mathcal{M}$-measurable:}
    
    Write $f = f^+ - f^-$ where $f^+ = \max(f, 0) \geq 0$ and $f^- = \max(-f, 0) \geq 0$.
    
    Define:
    \[
        h^+(x) = \Integral_Y f^+(x, y) \, d\nu(y), \qquad h^-(x) = \Integral_Y f^-(x, y) \, d\nu(y)
    \]
    
    By Tonelli, $h^+$ and $h^-$ are $\mathcal{M}$-measurable (as functions $X \to [0, \infty]$).
    
    For $x \in X \setminus N$: $h(x) = h^+(x) - h^-(x)$, where both terms are finite.
    
    For $x \in N$: $h(x) = 0$.
    
    Since $N \in \mathcal{M}$ (as $N = g^{-1}(\{\infty\})$ and $g$ is $\mathcal{M}$-measurable), 
    the function $h$ is $\mathcal{M}$-measurable.
    
    \textit{$h \in \Lp{1}(\mu)$:}
    
    For $x \in X \setminus N$:
    \[
        |h(x)| = \left|\Integral_Y f_x \, d\nu\right| \leq \Integral_Y |f_x| \, d\nu = g(x)
    \]
    
    For $x \in N$: $|h(x)| = 0 \leq g(x)$ (trivially, though $g(x) = \infty$ on $N$).
    
    Thus $|h| \leq g$ on all of $X$, and:
    \[
        \Integral_X |h| \, d\mu \leq \Integral_X g \, d\mu < \infty
    \]
    so $h \in \Lp{1}(\mu)$.
    
    \bigskip
    \textbf{Step 4: Compute the iterated integral using positive and negative parts.}
    
    Since $f = f^+ - f^-$ with $f^+, f^- \geq 0$, by Tonelli:
    \[
        \Integral_{X \times Y} f^+ \, d(\mu \otimes \nu) = \Integral_X h^+(x) \, d\mu(x)
    \]
    \[
        \Integral_{X \times Y} f^- \, d(\mu \otimes \nu) = \Integral_X h^-(x) \, d\mu(x)
    \]
    
    \textit{Both sides are finite:}
    \[
        \Integral_{X \times Y} f^+ \, d(\mu \otimes \nu) \leq \Integral_{X \times Y} |f| \, d(\mu \otimes \nu) < \infty
    \]
    and similarly for $f^-$.
    
    \textit{Subtracting:}
    Since both integrals are finite, we can subtract:
    \begin{align*}
        \Integral_{X \times Y} f \, d(\mu \otimes \nu) 
        &= \Integral_{X \times Y} f^+ \, d(\mu \otimes \nu) - \Integral_{X \times Y} f^- \, d(\mu \otimes \nu) \\
        &= \Integral_X h^+(x) \, d\mu - \Integral_X h^-(x) \, d\mu \\
        &= \Integral_X (h^+(x) - h^-(x)) \, d\mu \quad \text{(since both are in $\Lp{1}(\mu)$)}
    \end{align*}
    
    Now, on $X \setminus N$: $h^+(x) - h^-(x) = h(x) = \Integral_Y f_x \, d\nu$.
    
    On $N$: $\mu(N) = 0$, so the integrand's value there is irrelevant.
    
    Therefore:
    \[
        \Integral_{X \times Y} f \, d(\mu \otimes \nu) = \Integral_X h(x) \, d\mu = \Integral_{X \setminus N} \left(\Integral_Y f(x, y) \, d\nu(y)\right) d\mu(x)
    \]
    
    (The last equality uses that $h = 0$ on $N$ and $\mu(N) = 0$, so integrating over $X$ or $X \setminus N$ gives the same result.)
    
    By symmetry:
    \[
        \Integral_{X \times Y} f \, d(\mu \otimes \nu) = \Integral_{Y \setminus M} \left(\Integral_X f(x, y) \, d\mu(x)\right) d\nu(y)
    \]
    
    This completes the proof.
\end{proof}

\begin{corollary}
    If $f$ is $(\mathcal{M} \otimes \mathcal{N})$-measurable and one of the iterated integrals 
    $\Integral_X \left(\Integral_Y |f| \, d\nu\right) d\mu$ or $\Integral_Y \left(\Integral_X |f| \, d\mu\right) d\nu$ 
    is finite, then $f \in \Lp{1}(\mu \otimes \nu)$ and Fubini applies.
\end{corollary}

\begin{proof}
    By Tonelli, the iterated integral of $|f|$ equals $\Integral_{X \times Y} |f| \, d(\mu \otimes \nu)$.
    If this is finite, then $f \in \Lp{1}(\mu \otimes \nu)$.
\end{proof}

\begin{remark}[When Fubini Fails Without Integrability]
    The integrability condition in Fubini is essential. Consider $X = Y = [0, 1]$ with Lebesgue measure 
    and $f(x, y) = \frac{x^2 - y^2}{(x^2 + y^2)^2}$ for $(x, y) \neq (0, 0)$. Then the two iterated 
    integrals exist but are \textbf{not equal}:
    \[
        \Integral_{[0,1]} \left(\Integral_{[0,1]} f(x, y) \, d\lambda(y)\right) d\lambda(x) = \frac{\pi}{4}, \quad
        \Integral_{[0,1]} \left(\Integral_{[0,1]} f(x, y) \, d\lambda(x)\right) d\lambda(y) = -\frac{\pi}{4}
    \]
    This happens because $f \notin \Lp{1}(\lambda \otimes \lambda)$.
\end{remark}


\section{Consequences of Fubini-Tonelli}

\begin{intuition}
    Fubini-Tonelli is a powerful tool that enables us to derive several classical results 
    by viewing one-dimensional problems as sections of two-dimensional product spaces. 
    Two beautiful applications are:
    \begin{itemize}
        \item \textbf{Young's Inequality:} Derived by computing areas in the plane
        \item \textbf{Layer Cake Formula:} Expresses integrals via the distribution function
    \end{itemize}
\end{intuition}

\bigskip
\subsection{Young's Inequality}

\begin{theorem}[Young's Inequality --- Integral Form]
    Let $\varphi: [0, \infty) \to [0, \infty)$ be continuous, strictly increasing, with $\varphi(0) = 0$ 
    and $\lim_{t \to \infty} \varphi(t) = \infty$. Let $\psi = \varphi^{-1}$ be its inverse. Then:
    \[
        ab \leq \Integral_{[0,a]} \varphi(t) \, d\lambda(t) + \Integral_{[0,b]} \psi(s) \, d\lambda(s)
    \]
    with equality iff $b = \varphi(a)$.
\end{theorem}

\begin{proof}
    \textbf{Step 1: Set up the regions.}
    
    Consider the rectangle $R = [0, a] \times [0, b] \subseteq \bb{R}^2$. Define:
    \begin{align*}
        E_1 &= \{ (t, s) \in \bb{R}^2 \mid 0 \leq t \leq a, \, 0 \leq s \leq \varphi(t) \} \\
        E_2 &= \{ (t, s) \in \bb{R}^2 \mid 0 \leq s \leq b, \, 0 \leq t \leq \psi(s) \}
    \end{align*}
    
    \medskip
    \textbf{Step 2: Show $R \subseteq E_1 \cup E_2$.}
    
    Let $(t, s) \in R$. If $s \leq \varphi(t)$, then $(t,s) \in E_1$. 
    If $s > \varphi(t)$, then $t < \psi(s)$ (since $\varphi$ is strictly increasing), so $(t,s) \in E_2$.
    
    \medskip
    \textbf{Step 3: Compute the measures using Tonelli.}
    
    By Tonelli's theorem:
    \[
        \lambda^2(E_1) = \Integral_{[0,a]} \varphi(t) \, d\lambda(t), \quad 
        \lambda^2(E_2) = \Integral_{[0,b]} \psi(s) \, d\lambda(s)
    \]
    
    \medskip
    \textbf{Step 4: Conclude.}
    
    Since $R \subseteq E_1 \cup E_2$:
    \[
        ab = \lambda^2(R) \leq \lambda^2(E_1 \cup E_2) \leq \lambda^2(E_1) + \lambda^2(E_2) 
        = \Integral_{[0,a]} \varphi(t) \, d\lambda(t) + \Integral_{[0,b]} \psi(s) \, d\lambda(s)
    \]
    
    \medskip
    \textbf{Step 5: Equality case.}
    
    Equality holds iff $R = E_1 \cup E_2$ and $E_1 \cap E_2$ has measure zero.
    This occurs iff $b = \varphi(a)$, since then the point $(a, b)$ lies on the graph of $\varphi$, 
    and $E_1, E_2$ partition $R$ (up to the graph, which has measure zero).
\end{proof}

\begin{corollary}[Young's Inequality for $p$-powers]
    Let $1 < p, q < \infty$ with $\frac{1}{p} + \frac{1}{q} = 1$. Then for all $a, b \geq 0$:
    \[
        ab \leq \frac{a^p}{p} + \frac{b^q}{q}
    \]
    with equality iff $a^p = b^q$.
\end{corollary}

\begin{proof}
    Apply the integral form with $\varphi(t) = t^{p-1}$, so $\psi(s) = s^{q-1}$ 
    (since $(p-1)(q-1) = 1$ follows from $\frac{1}{p} + \frac{1}{q} = 1$).
    
    Compute $\Integral_{[0,a]} t^{p-1} \, d\lambda(t) = \frac{a^p}{p}$ and $\Integral_{[0,b]} s^{q-1} \, d\lambda(s) = \frac{b^q}{q}$.
    
    Equality holds iff $b = \varphi(a) = a^{p-1}$, i.e., $b^q = a^{(p-1)q} = a^p$.
\end{proof}

\bigskip
\subsection{Layer Cake Formula}

\begin{theorem}[Layer Cake Representation]
    Let $(X, \mathcal{M}, \mu)$ be a measure space and $f: X \to [0, \infty]$ be measurable. Then:
    \[
        \Integral_X f \, d\mu = \Integral_{[0,\infty)} \mu(\{f > t\}) \, d\lambda(t)
    \]
    where $\lambda$ is Lebesgue measure on $[0, \infty)$.
\end{theorem}

\begin{intuition}
    The layer cake formula says we can compute the integral of a non-negative function 
    by ``stacking horizontal slices''. For each height $t \geq 0$, we measure the set 
    $\{f > t\}$ (where the function exceeds $t$), and integrate these measures.
    
    Geometrically, if we think of $f$ as defining a ``cake'' with height $f(x)$ at each 
    point $x$, then instead of integrating vertically (summing $f(x) \cdot d\mu(x)$), 
    we integrate horizontally (summing $\mu(\{f > t\}) \cdot dt$).
    
    This is Cavalieri's principle for measure spaces!
\end{intuition}

\begin{proof}
    We apply Tonelli's theorem on the product space $X \times [0, \infty)$ with measures 
    $\mu$ and $\lambda$.
    
    \bigskip
    \textbf{Step 1: Define the subgraph.}
    
    Let $E = \{(x, t) \in X \times [0, \infty) \mid 0 \leq t < f(x)\}$ be the ``subgraph'' of $f$.
    
    \medskip
    \textbf{Step 2: Show $E$ is measurable.}
    
    Define $g: X \times [0, \infty) \to \bb{R}$ by $g(x, t) = f(x) - t$. Since $f$ is measurable 
    and projection onto the second coordinate is measurable:
    \[
        E = g^{-1}((0, \infty)) = \{(x, t) \mid f(x) > t\} \in \mathcal{M} \otimes \BorelAlg{[0, \infty)}
    \]
    
    \medskip
    \textbf{Step 3: Compute $(\mu \otimes \lambda)(E)$ via $x$-sections (Tonelli).}
    
    For fixed $x \in X$, the $x$-section is:
    \[
        E_x = \{t \in [0, \infty) \mid t < f(x)\} = [0, f(x))
    \]
    Thus $\lambda(E_x) = f(x)$ (with convention $\lambda([0, \infty)) = \infty$).
    
    By Tonelli:
    \[
        (\mu \otimes \lambda)(E) = \Integral_X \lambda(E_x) \, d\mu(x) = \Integral_X f(x) \, d\mu(x)
    \]
    
    \medskip
    \textbf{Step 4: Compute $(\mu \otimes \lambda)(E)$ via $t$-sections (Tonelli).}
    
    For fixed $t \in [0, \infty)$, the $t$-section is:
    \[
        E^t = \{x \in X \mid f(x) > t\} = \{f > t\}
    \]
    Thus $\mu(E^t) = \mu(\{f > t\})$.
    
    By Tonelli:
    \[
        (\mu \otimes \lambda)(E) = \Integral_{[0,\infty)} \mu(E^t) \, d\lambda(t) = \Integral_{[0,\infty)} \mu(\{f > t\}) \, d\lambda(t)
    \]
    
    \medskip
    \textbf{Step 5: Conclude.}
    
    Both expressions equal $(\mu \otimes \lambda)(E)$, so:
    \[
        \Integral_X f \, d\mu = \Integral_{[0,\infty)} \mu(\{f > t\}) \, d\lambda(t)
    \]
\end{proof}

\begin{remark}
    The formula also holds with $\{f \geq t\}$ instead of $\{f > t\}$:
    \[
        \Integral_X f \, d\mu = \Integral_{[0,\infty)} \mu(\{f \geq t\}) \, d\lambda(t)
    \]
    This is because $\{f > t\}$ and $\{f \geq t\}$ differ only on $\{f = t\}$, which forms 
    a family of disjoint sets as $t$ varies, hence has $\lambda$-measure zero in the integral.
\end{remark}

\begin{corollary}[Power Formula]
    For $f \geq 0$ measurable and $p > 0$:
    \[
        \Integral_X f^p \, d\mu = p \Integral_{[0,\infty)} t^{p-1} \mu(\{f > t\}) \, d\lambda(t)
    \]
\end{corollary}

\begin{proof}
    Apply the layer cake formula to $f^p$:
    \[
        \Integral_X f^p \, d\mu = \Integral_{[0,\infty)} \mu(\{f^p > s\}) \, d\lambda(s)
    \]
    Substitute $s = t^p$, so $ds = p t^{p-1} dt$:
    \[
        = \Integral_{[0,\infty)} \mu(\{f > t\}) \cdot p t^{p-1} \, d\lambda(t) = p \Integral_{[0,\infty)} t^{p-1} \mu(\{f > t\}) \, d\lambda(t)
    \]
\end{proof}

\bigskip
\begin{compendium}

\textbf{Product $\sigma$-Algebra:} For measurable spaces $(X, \mathcal{M})$ and $(Y, \mathcal{N})$:
\[
    \mathcal{M} \otimes \mathcal{N} = \GenSigmaAlg{\{A \times B \mid A \in \mathcal{M}, B \in \mathcal{N}\}}
\]

\textbf{Borel $\sigma$-Algebras:}
\begin{itemize}
    \item General: $\BorelAlg{X} \otimes \BorelAlg{Y} \subseteq \BorelAlg{X \times Y}$
    \item Second countable: $\BorelAlg{X} \otimes \BorelAlg{Y} = \BorelAlg{X \times Y}$
\end{itemize}

\textbf{Sections:} For $E \subseteq X \times Y$:
\begin{itemize}
    \item $x$-section: $E_x = \{y \in Y \mid (x, y) \in E\}$
    \item $y$-section: $E^y = \{x \in X \mid (x, y) \in E\}$
\end{itemize}
If $E \in \mathcal{M} \otimes \mathcal{N}$, then $E_x \in \mathcal{N}$ and $E^y \in \mathcal{M}$.

\textbf{Product Measure:} For $\sigma$-finite $(X, \mathcal{M}, \mu)$ and $(Y, \mathcal{N}, \nu)$:
\[
    (\mu \otimes \nu)(E) = \Integral_X \nu(E_x) \, d\mu(x) = \Integral_Y \mu(E^y) \, d\nu(y)
\]
Uniquely characterized by $(\mu \otimes \nu)(A \times B) = \mu(A) \cdot \nu(B)$.

\textbf{Tonelli (Non-negative):} For $f \geq 0$ measurable:
\[
    \Integral_{X \times Y} f \, d(\mu \otimes \nu) = \Integral_X \left(\Integral_Y f(x,y) \, d\nu\right) d\mu = \Integral_Y \left(\Integral_X f(x,y) \, d\mu\right) d\nu
\]
No integrability assumption needed.

\textbf{Fubini (Integrable):} For $f \in \Lp{1}(\mu \otimes \nu)$: same equality holds, with iterated integrals 
over $X \setminus N$ and $Y \setminus M$ where $N, M$ are null sets.

\textbf{Completeness Issue:} Product of complete spaces need not be complete. Sections may fail to be 
measurable on null sets, necessitating integration over $X \setminus N$.

\textbf{Young's Inequality (Fubini Consequence):}
\[
    ab \leq \Integral_{[0,a]} \varphi(t) \, d\lambda(t) + \Integral_{[0,b]} \psi(s) \, d\lambda(s)
\]
where $\psi = \varphi^{-1}$. For $p$-powers: $ab \leq \frac{a^p}{p} + \frac{b^q}{q}$ with $\frac{1}{p} + \frac{1}{q} = 1$.

\textbf{Layer Cake Formula:} For $f \geq 0$ measurable:
\[
    \Integral_X f \, d\mu = \Integral_{[0,\infty)} \mu(\{f > t\}) \, d\lambda(t)
\]

\textbf{Power Formula:} For $f \geq 0$ and $p > 0$:
\[
    \Integral_X f^p \, d\mu = p \Integral_{[0,\infty)} t^{p-1} \mu(\{f > t\}) \, d\lambda(t)
\]

\end{compendium}

